% ============================================================================
% AI Agentic Sandbox -- paper draft
% ============================================================================
% Purpose: a discussion draft for Mike + PI, shared via Overleaf. Section/
% subsection structure follows the outline agreed 2026-08-25; body text below
% is now real prose, filled in from docs/vision/vision_paper.tex, the current
% agent-sandbox/ codebase (architecture.md, README.md, future-work-roadmap.md,
% and the schema/test source itself), and resource-scheduler/README.md.
% Remaining \todo{...} markers are genuine open items -- author list, figures,
% a real Related Work literature search, and the resource-scheduler
% co-worker's in-progress port/evaluation.
%
% Focus, per PI discussion (2026-08-25): the Sandbox itself -- what it is,
% how it currently functions, and the merit of the idea -- is the paper.
% agentic-ml-classification and resource-scheduler are supporting case
% studies / validation evidence, not the subject of the paper themselves.
%
% Format note: using IEEEtran (10pt, conference) because it is the required
% template for 4 of the 6 venues under consideration (SANER, ICSE NIER,
% ICSE SEIP, and likely SEAMS). If the target venue turns out to be FSE
% (ACM acmart) or the EMSE journal (Springer template), the section
% structure below carries over directly -- only the \documentclass and
% front matter need to change. See docs/paper_outline.md for the full
% venue-by-venue breakdown this draft was scoped from.
% ============================================================================

\documentclass[10pt,conference]{IEEEtran}

\usepackage{cite}
\usepackage{amsmath,amssymb}
\usepackage{graphicx}
\usepackage{textcomp}
\usepackage{xcolor}
\usepackage{booktabs}
\usepackage{hyperref}

\newcommand{\todo}[1]{\textcolor{red}{\textbf{[TODO: #1]}}}
\newcommand{\scaffold}[1]{\textit{#1}}

\begin{document}

\title{The AI Agentic Sandbox: An Innovation Lab for Trustworthy,\\
Bring-Your-Own-Agent Multi-Agent Systems in Manufacturing}

\author{
\IEEEauthorblockN{\todo{Author list -- vision paper still says ``Authors: to be added''; resolve with PI before any submission, including this draft's authorship order}}
\IEEEauthorblockA{Rochester Institute of Technology, Rochester, NY, USA}
}

\maketitle

% ----------------------------------------------------------------------------
\begin{abstract}
Multi-agent systems (MAS) built from large language model (LLM) agents are an
attractive fit for manufacturing's distributed, divided-labor structure, but
their design, optimization, and evaluation remain heuristic-driven, and
university-industry collaboration on them today is transactional and
one-off -- an upfront IP/licensing agreement before any research happens,
which does not scale to an era where every partner is independently building
agents. We present the AI Agentic Sandbox, a MAS infrastructure at RIT built
on a schema-first foundation (declarative \texttt{AgentSpec}/\texttt{RunSpec}
specifications), MCP-mediated tool access with per-binding allowlists, an
append-only observable event log, and a default-deny network isolation
boundary, that lets partners bring their own agents (BYOA) and swap or
compare them inside shared workflows without re-architecting the workflow
around them. We report precisely what is built and tested today --
declarative agent/pipeline authoring, a deterministic multi-step pipeline
runtime with agent and gate steps, a pre-run operator layer
(\texttt{SWAP\_AGENT}/\texttt{ALTER\_MODEL}/\texttt{ALTER\_WORKFLOW}) with a
forensic audit trail, and a service/UI layer, together covered by 160
automated tests -- versus what remains proposed: mid-run operation
mutation, LLM-driven dynamic sub-agent spawning, multi-tenant per-user
isolation, and cross-institution MCP governance. As validating evidence, we
summarize two independently-built MAS applications -- a six-agent ML
classification pipeline and a direct agent-to-agent resource scheduler --
that converge, without having shared a codebase, on the same underlying
principle: an LLM agent's output is only ever a proposal that a
deterministic layer independently re-validates before it affects anything.
That convergence is offered as evidence the Sandbox's core trust-boundary
pattern generalizes, not as a claim that either system is the Sandbox. We
close with a five-part research agenda -- extended MCP governance, automated
MAS harness engineering, MAS-specific cybersecurity, system-level MAS
evaluation, and an agentic-marketplace alternative to traditional tech
transfer -- that the Sandbox's build has surfaced as open problems rather
than solved ones.
\end{abstract}

\begin{IEEEkeywords}
agentic AI, multi-agent systems, trust boundary, manufacturing, Model
Context Protocol, university-industry collaboration
\end{IEEEkeywords}

% ============================================================================
\section{Introduction}
\label{sec:intro}
% ----------------------------------------------------------------------------
AI agents, functioning as digital workers, have become an integral part of
enterprise systems. Unlike conventional chat assistants, these systems are
increasingly given direct access to code, filesystems, shells, browsers, and
external services. IDC projects more than 1 billion actively deployed AI
agents by 2029, executing roughly 217 billion actions a day, and forecasts
agentic AI will exceed 26\% of worldwide IT spending -- \$1.3 trillion -- that
same year~\cite{perilli2026agentic}. Multi-agent systems (MAS) extend the
capability of a single specialized agent by composing multiple agents that
coordinate or compete toward a shared objective: distributed rather than
centralized, divided among specialized agents, and capable of exhibiting
behavior that emerges from collective interaction rather than any single
agent's design~\cite{Barbosa2010,Leitao2012,Bussmann2004}. Manufacturing is
an attractive domain for this paradigm -- IIoT sensors, edge/cloud computing,
and digital twins are already reshaping it from reactive and scheduled
operations toward predictive, resilient, real-time decision-making -- but the
design, optimization, and evaluation of MAS applications remain
heuristic-driven, with no standardized design patterns or evaluation
schemes, and LLM-powered MAS specifically remain relatively unexplored in
industrial contexts~\cite{Farahani_2026}.

A second, largely orthogonal gap sits alongside the technical one:
university-industry MAS collaboration today is transactional and one-off.
Achieving upfront agreement on deliverables, IP sharing, and licensing terms
is, in practice, often the actual blocker to collaboration -- not a shortage
of good ideas~\cite{hargadon2003breakthroughs}. That model does not scale to
an era where every partner is independently building agents and wants to
experiment with combining them before any of it is worth negotiating a
license over.

This paper describes our answer to both gaps together: the AI Agentic
Sandbox, a shared, safe infrastructure where university and industry
partners bring their own agents, swap them into shared manufacturing
workflows, and evaluate system-level behavior before any of it touches
production. The Sandbox's value proposition is infrastructural, not a
specific manufacturing application -- the paper's two case studies exist to
validate that the underlying trust-boundary pattern generalizes across
domains and orchestration topologies, not to present either as a finished
product.

\textbf{Contributions.}
\begin{itemize}
  \item A schema-first MAS infrastructure (\texttt{AgentSpec}/\texttt{RunSpec}/
  \texttt{PipelineSpec}) that makes agent composition, tool access, and
  observability declarative rather than bespoke orchestration code
  (\S\ref{sec:infra}).
  \item A concrete, implemented, and tested Bring-Your-Own-Agent mechanism --
  \texttt{SWAP\_AGENT}, \texttt{ALTER\_MODEL}, \texttt{ALTER\_WORKFLOW} as
  validated, re-validated, and forensically logged pre-run spec mutations --
  together with two independently-built MAS applications used as evidence
  the underlying design pattern generalizes (\S\ref{sec:infra},
  \S\ref{sec:casestudies}).
  \item An honest, code-verified implementation-status account: 160
  automated tests currently pass across the Sandbox's three packages,
  distinguishing what that coverage actually covers from what remains
  proposed -- mid-run operation mutation, LLM-driven dynamic sub-agent
  spawning, and multi-tenant isolation chief among them
  (\S\ref{sec:infra}, \S\ref{sec:limitations}).
  \item A research agenda for the open problems the Sandbox's build
  surfaces: cross-institution MCP governance, automated MAS harness
  engineering, and MAS-specific cybersecurity threat modeling
  (\S\ref{sec:merit}).
\end{itemize}

% ============================================================================
\section{Motivation and Background}
\label{sec:motivation}
% ----------------------------------------------------------------------------

\subsection{Agentic Systems}
We adopt NIST's definition of an agent: software that interacts with its
environment, receives information, and undertakes self-directed action
toward an externally-specified goal~\cite{nist_glossary_agent}. Concretely,
an agent combines three elements -- a model, a set of tools, and a loop that
runs until something decides to stop. A large language model powers the
reasoning and planning that let the agent generate actions, with or without
human supervision, toward a specific goal; tools are reached almost
exclusively through Model Context Protocol (MCP) servers, an open standard
that provides one consistent interface between a model and everything
outside it. Unlike a rigid rule-based system, an LLM-driven agent adapts,
learns, and responds dynamically to situations its designer did not
enumerate in advance -- but that same generative flexibility is a liability
when the desired output has an objectively correct answer. The pairing this
whole paper depends on is: the generative power of a probabilistic LLM,
paired with the logical certainty of a symbolic \emph{harness} -- domain
rules, hard constraints, and human-operator validation or guards. A
probabilistic LLM output is a property of the model, not a defect in it;
research consistently shows that choosing the model is not enough --
orchestration harness, memory strategy, context handling, validation path,
and checkpointing design drive practical agent behavior far more than model
selection alone. \todo{Reuse Fig.~2 from the vision PDF (agent/environment
interface diagram) or redraw it.}

\subsection{Multi-Agent Systems}
When multiple agents coordinate or compete in a shared environment toward a
complex, shared objective, they form an agentic \emph{system}. An
application built this way is a workflow composed of an arbitrary set of
specialized agents. Two topologies recur throughout this paper's case
studies (\S\ref{sec:casestudies}), not as a binary choice but as two ends of
a spectrum: an \emph{orchestrated} MAS, where a central coordinator sequences
every agent's turn, and a \emph{peer-to-peer} MAS, where agents negotiate
directly, agent-to-agent (A2A), for the parts of a workflow that are
naturally negotiation-shaped. Both remain compatible with the same
underlying trust-boundary discipline this paper argues for -- the
orchestration topology and the propose/decide authority split are
independent design axes, not the same decision.

\subsection{Manufacturing as the Driving Domain}
Interconnected fleets of cyber-physical systems, IIoT sensors, advanced
AI/ML, edge and cloud computing, and digital twin technologies are reshaping
manufacturing from reactive and scheduled operations toward predictive,
resilient, real-time decision-making ecosystems. The MAS paradigm is an
attractive fit for manufacturing specifically because manufacturing systems
already share MAS's own structural attributes: a distributed nature
(overall functioning achieved through interaction among distributed,
cooperative units), a division of labor (among specialized stations or
processes), and emergence from collective, simple behavior (complex
plant-level outcomes aggregated from many local decisions)
\cite{Barbosa2010,Leitao2012,Bussmann2004}. \todo{cite vision paper's
manufacturing-systems references (Farahani et al., Liao et al., etc.) more
completely once Related Work (\S\ref{sec:related}) is drafted.}

% ============================================================================
\section{The AI Agentic Sandbox}
\label{sec:sandbox}
% ----------------------------------------------------------------------------
\emph{The AI Agentic Sandbox is a MAS infrastructure where universities and
industry partners can collaboratively develop, experiment with, and evaluate
multi-agent AI applications to advance innovation in manufacturing systems.}
The Sandbox offers enterprises a \emph{safe space} to: (i) co-innovate with
partners; (ii) experiment with and evaluate how autonomous agents behave
within workflows; and (iii) optimize agentic applications before they are
trusted with production stakes -- an agentic innovation lab built as
technology infrastructure, addressing those needs specifically for
manufacturing-system workflows.

\subsection{Goals}
\label{sec:goals}
\begin{description}
  \item[Goal 1 -- BYOA] Every partner can deploy proprietary agents and
  multi-agentic workflows to experiment against other agents already in
  the Sandbox.
  \item[Goal 2 -- Holistic evaluation] What-if analysis and architecture
  trade-off evaluation for specific MAS workflows, so partners understand
  emergent system behavior and identify risk before deployment.
  \item[Goal 3 -- Translation avenue] University-based AI research results,
  implemented as specialized agents and workflows, get a path to industry
  adoption via an \emph{agentic marketplace} rather than the traditional
  one-off tech-transfer/licensing pipeline.
  \item[Goal 4 -- RIT value-add] Cyber risk assessment/mitigation,
  explainable agency, and safe-AI-in-MAS-context as capabilities RIT
  specifically contributes.
\end{description}

\subsection{Sandbox Design}
\label{sec:design}
The Sandbox's infrastructure rests on three foundations, each realized (or
partially realized -- see \S\ref{sec:infra}) in the actual codebase described
below:
\begin{itemize}
  \item \textbf{Partners} -- university, industry, and government entities,
  each a team of humans \emph{and} agents, potentially operating under
  different institutional priorities (fundamental research vs. market
  outcomes vs. policy/societal impact). A university research faculty
  member may prioritize scholarship where an industry partner prioritizes
  market-driven outcomes; the Sandbox is designed to let both work inside
  the same infrastructure without either priority overriding the other.
  \item \textbf{Agent Store} -- a curated set of specialized agents and MAS
  workflows, contributed by partners or drawn from open sources, seeded
  with foundational manufacturing applications (\S\ref{sec:seedapps}). In
  the current implementation this is realized concretely as a directory of
  versioned \texttt{AgentSpec}/\texttt{PipelineSpec} YAML files, each
  addressable and swappable independently (\S\ref{sec:infra}).
  \item \textbf{Agentic Operations} -- utilities for building new
  applications by adapting existing workflows. Three are implemented today,
  as validated pre-run spec mutations rather than raw file edits:
  \texttt{SWAP\_AGENT} (repoint one pipeline step at a different agent
  without touching the rest of the workflow), \texttt{ALTER\_MODEL}
  (replace an agent's model configuration wholesale, to experiment with an
  alternate LLM), and \texttt{ALTER\_WORKFLOW} (replace a pipeline's step
  list wholesale -- the general-purpose edit of which
  \texttt{SWAP\_AGENT} is a common, narrower special case). Each operation
  is independently re-validated against the current spec before it is
  applied and produces an append-only, before/after/diff forensic record
  (\S\ref{sec:infra}). \todo{These three operators are pre-run only today
  (they mutate a spec sitting in a specs directory, not a live run's
  in-memory state); \emph{mid-run} operation mutation and further operator
  types (e.g.\ \texttt{SPAWN\_SUBAGENT}) remain future work -- state this
  scope boundary plainly rather than implying the operator list is
  closed.}
\end{itemize}

\todo{Figure: reuse/redraw vision-PDF Fig.~1 (``AI Agentic Sandbox
Vision'') -- Partner A/B, agent store, orchestration, SWAP\_AGENT /
ALTER\_MODEL / ALTER\_WORKFLOW callouts.}

\subsection{Seed Applications}
\label{sec:seedapps}
The Sandbox is seeded with three foundational manufacturing application
domains, chosen to cover a representative array of manufacturing MAS
scenarios rather than an exhaustive one: \emph{predictive maintenance},
which forecasts machine failures before they occur, enabling a more adaptive
production system; \emph{planning and scheduling}, which generates
up-to-date, optimized production schedules by jointly considering demand and
shop-floor resource availability; and \emph{supply chains and virtual
organization}, which demand the frequent re-formation of partnerships to
sustain optimization and responsiveness under changing conditions. This
paper's two case studies (\S\ref{sec:casestudies}) concretely realize the
first two of these three -- predictive maintenance and resource scheduling,
respectively; the supply-chain/virtual-organization seed application remains
unbuilt (\S\ref{sec:limitations}).

% ============================================================================
\section{Sandbox Infrastructure}
\label{sec:infra}
% ----------------------------------------------------------------------------
This section describes how the Sandbox actually works today, as code, and is
precise throughout about built-vs-proposed; see \S\ref{sec:status} for a
consolidated accounting.

\subsection{Schema-First Foundation: AgentSpec, RunSpec, and PipelineSpec}
Every capability an agent can reach is registered through one interface
(MCP); every action is a typed, replayable event; every agent and every
workflow is described by a declarative spec rather than bespoke
orchestration code -- the same schema-first approach independently converged
on by Meta's Agents Research Environments
(ARE)~\cite{froger2025arescalingagentenvironments}, which structures agent
evaluation around typed environments, apps, and events rather than ad hoc
harness code. An \texttt{AgentSpec} declares an agent's identity, system
prompt, model configuration, bound MCP servers (each with a transport, an
explicit tool allowlist, and a redaction policy), and any sub-agents it may
delegate to; a \texttt{RunSpec} declares a single task against one
\texttt{AgentSpec}; a \texttt{PipelineSpec} declares a sequence of steps --
either an ordinary agent step or a deterministic \emph{gate} step
(\S\ref{sec:pipelines}) -- chaining multiple agents into one
harness-supervised workflow. Because an agent's tools, model, and workflow
position are \emph{data}, not code, a partner's proprietary agent can
replace a seed agent by swapping one binding, without touching the
surrounding workflow -- the concrete mechanism behind the
\texttt{SWAP\_AGENT}/\texttt{ALTER\_MODEL} operators from \S\ref{sec:design}.
All of these schemas are pydantic v2 models, and are additionally exported
as versioned JSON Schema documents so a non-Python consumer -- a future
hosted UI, or another language's agent implementation -- has a stable
contract to code against without importing the Python package directly.

\subsection{Harness: Strands Agents SDK}
\label{sec:strands}
The actual tool-calling loop, model retries, and MCP session lifecycle are
not hand-rolled in the Sandbox's own code -- they are delegated to
\href{https://strandsagents.com}{Strands Agents}, an open-source agent SDK.
sandbox-core's own code is the YAML/pydantic authoring surface plus a thin
adapter (\texttt{strands\_adapter.py}) that turns a validated
\texttt{AgentSpec} into a \texttt{strands.Agent}: it resolves the model API
key, builds an OpenAI-compatible model client against \texttt{ModelConfig}'s
endpoint (any \texttt{/chat/completions}-compatible target -- RIT's gateway,
a local vLLM instance, or a partner's own endpoint), and constructs one MCP
client per \texttt{McpServerBinding} with the binding's tool allowlist
enforced before the agent ever sees a tool's schema. What the Sandbox's own
code still owns, and had to preserve exactly across this migration, is
everything downstream of Strands' hooks: a \texttt{SandboxEventHooks}
provider translates Strands' own before/after model- and tool-call events
into the Sandbox's typed \texttt{Event} union, applying each binding's
redaction policy (full / hashed / metadata-only) before persisting a tool
result -- the model itself always receives the unredacted result; only what
lands on disk is policy-filtered. \texttt{execute\_run()}'s signature and
behavior toward its callers (the CLI, the server) were kept unchanged
through the migration, so nothing above this layer needed to change to
benefit from it. What Strands provided "for free": correct MCP session
lifecycle management, per-call retry behavior, and turn-limit enforcement
that a hand-rolled loop previously had to reimplement and re-verify itself.

\subsection{MCP-Mediated Tool Surface and Access Control}
Agents never call tools directly; every capability is exposed through an MCP
server binding with an explicit allowlist, enforced via Strands' tool
filters before a call reaches the server -- a filtered-out tool is never
loaded onto the agent in the first place, not merely hidden after the fact.
Credentials resolve through a dedicated credential layer (a flat, gitignored
YAML store today, resolved by reference) -- a reference-only handoff that
keeps raw secrets out of an agent's context entirely; the only place a
credential's real value ever exists is inside the resolver call itself, and
there is deliberately no API route that can read one back out once written.
This is the concrete mechanism this paper proposes extending toward
institution-specific, policy-driven access control (\S\ref{sec:merit}).

\subsection{Observability: the Event Log}
Every model call, tool call, and terminal outcome is written to an
append-only, per-run typed event log (one JSON line per event, with a
monotonic per-run sequence number), with the same configurable redaction
policy from \S\ref{sec:strands} applied per binding before persistence. This
is the substrate for explainable agency (Goal 4, \S\ref{sec:goals}): because
every agent's actions are captured uniformly regardless of model or tools
used, a partner can audit or replay a run -- including one served live over
SSE to a browser mid-run, or replayed after the fact straight from disk --
across agents using different LLMs.

\subsection{Isolation Boundary}
Every MCP server registered in the Sandbox's agent store is local to the
Sandbox; by default, no agent has outbound access to external networks,
databases, or unregistered MCP servers. This is what makes the Sandbox a
genuine "safe space" structurally, not just by policy convention -- enforced
at the same allowlist boundary that mediates every tool call
(\S\ref{sec:strands}), not left to an agent's good behavior or a partner's
configuration.

\subsection{Deterministic Multi-Agent Pipelines}
\label{sec:pipelines}
Multi-agent composition in the Sandbox today is \emph{harness-driven}, not
LLM-driven: a \texttt{PipelineSpec} is a sequence of steps the harness walks
deterministically, rather than one agent autonomously deciding to spawn
another. Two step kinds compose: an ordinary \texttt{PipelineStep} is a
complete agent run in its own right, with its task built from a template
that can reference the pipeline's seed task or a prior step's output; a
\texttt{GateStep} is a deterministic, non-LLM Python callable -- no agent run
backs it -- that inspects every prior step's output and returns a decision
string, which the pipeline's declared routing table maps to the next
step, including \emph{backward}, expressing a reject$\to$retry loop a
strictly linear step list cannot. This is the "agents propose, harness
decides" pattern from \S\ref{sec:casestudy1} made generic and reusable at
the infrastructure level, not something particular to the ML-classification
case study: a gate step is deliberately not sandboxed against untrusted code
(\S\ref{sec:limitations}), but it lets an operator's own deterministic
validation code -- reused unmodified from an existing package -- gate
whether a pipeline proceeds, exactly as \S\ref{sec:casestudy1}'s harness
gates a modeling candidate. A \texttt{max\_steps} cap (default 50) bounds
total step executions per run, the pipeline-level analogue of an
individual agent's own turn limit.

\subsection{Implementation Status}
\label{sec:status}
State this plainly and specifically, verified directly against the code and
test suite as of 2026-09-01 rather than against any internal milestone
label:
\begin{itemize}
  \item \textbf{Built and tested (160 automated tests, all passing):} the
  full schema layer (agent, run, event, pipeline, and operation
  specifications); a headless single-agent runtime built on Strands
  (\S\ref{sec:strands}) with MCP allowlist enforcement, credential
  resolution, and an append-only event log, CLI-drivable with no server
  running at all; a deterministic multi-agent pipeline runtime
  (\S\ref{sec:pipelines}) supporting both linear and reject/retry-loop
  workflows; a pre-run operator layer (\S\ref{sec:design}) --
  \texttt{SWAP\_AGENT}, \texttt{ALTER\_MODEL}, \texttt{ALTER\_WORKFLOW} --
  each independently re-validated and forensically logged before/after/diff
  on apply; and a full service and UI layer (a FastAPI service exposing
  agents, pipelines, credentials, and runs over REST plus live SSE
  streaming, and a React/Vite single-page app for authoring and watching
  runs). Test coverage splits 102 tests in the core library, 53 in the
  service layer, and 5 in the UI.
  \item \textbf{Not yet built:} \emph{mid-run} operation mutation (today's
  operators target a draft spec, never a live run's in-memory state);
  LLM-driven dynamic sub-agent spawning (\texttt{SubAgentBinding} and
  \texttt{AgentSpawnEvent} exist in the schema layer, but the runtime does
  not yet construct child agents from them -- deliberately not pursued as
  the primary multi-agent mechanism, since LLM-driven sequencing is the
  opposite of \S\ref{sec:pipelines}'s harness-driven design); sandboxed
  execution of \emph{untrusted} gate code (today's gate steps run trusted,
  in-process Python -- fine for an operator's own deterministic validation
  logic, not for arbitrary LLM-generated code, which needs AST-checked,
  subprocess-isolated execution as a separate, not-yet-built layer);
  multi-tenant per-user isolation (agents, pipelines, and runs are all
  global, on-disk, and shared today -- no concept of a user account
  exists yet); and an MCP gateway enforcing cross-institution governance
  policy (\S\ref{sec:merit}).
\end{itemize}

% ============================================================================
\section{Validating the Design Pattern: Case Studies}
\label{sec:casestudies}
% ----------------------------------------------------------------------------
These two projects are supporting evidence that the Sandbox's core design
pattern generalizes, not the paper's subject; this section stays
proportionate to that role rather than fully re-describing either system (see
\texttt{agentic-ml-classification/docs/ml-classification.tex} and
\texttt{resource-scheduler/README.md} for the full descriptions). Both were
built independently of the Sandbox codebase, before Sandbox infrastructure
existed to run them on, yet both converge on the same underlying principle:
an LLM agent's output is only ever a \emph{proposal} that a deterministic
layer independently re-validates before it affects anything. That
convergence, arrived at twice independently, is the evidence being offered
here, not a claim that either system \emph{is} the Sandbox.

\subsection{Case Study I: Multi-Agentic Predictive Maintenance}
\label{sec:casestudy1}
Six specialized agents -- intake, feature engineering, profiler, modeling,
verification, deep-dive -- operate under a deterministic harness holding
exclusive authority over splitting, metrics, leakage detection, and
sandboxed execution (Table~\ref{tab:authority}). Every agent output is a
JSON proposal the harness independently re-validates before acting on it. A
dynamic-orchestrator variant replaces the fixed six-step sequence with a
nine-entry agent catalog and an eighth agent, a planner, itself subject to
the same propose-then-verify discipline applied to control flow rather than
only content.

\begin{table}[h]
\caption{Division of authority between agents and harness (Case Study I).}
\label{tab:authority}
\centering
\small
\begin{tabular}{p{3.6cm}p{3.6cm}}
\toprule
\textbf{Agents propose} & \textbf{Harness decides} \\
\midrule
Target column and task framing & Train/val/test splits (group- and time-aware) \\
Feature drops and stateless derived features & Imputation, scaling, all fitted statistics \\
Model template and hyperparameters & Metric computation with bootstrap CIs \\
Plain-language narration and audit verdicts & Leakage gates, sandboxing, candidate promotion \\
Which catalog agent runs next & Precondition and identity validation before execution \\
\bottomrule
\end{tabular}
\end{table}

\begin{itemize}
  \item \textbf{Validated on:} NGAFID aviation predictive-maintenance data
  (22 sensor channels, \textasciitilde1\,Hz, 4.2\,GB source, streaming
  bounded-memory featurization, entity-aware splitting by aircraft), plus
  Titanic/Iris for binary/multiclass generalization.
  \item \textbf{Relevant finding for this paper specifically:} real-LLM
  evaluation (not just hermetic stubbed-client tests) surfaced two genuine
  findings a purely architectural description would not -- (1) the dynamic
  orchestrator matched or exceeded the fixed baseline (e.g., autonomously
  retrying a modeling step past a hardcoded candidate budget the baseline
  could not get past), and (2) the harness's state bookkeeping held under
  adversarial pressure: a prompt-injection string embedded in training
  data could not make the system assert a model was finalized unless
  finalization genuinely executed.
  \item \textbf{What this validates about the Sandbox pattern:} a MAS can
  give agents genuine judgment authority (target selection, feature
  proposals, template choice, even control-flow sequencing) while a
  non-agentic layer retains sole authority over everything
  correctness-critical -- the same authority split \S\ref{sec:infra}
  implements at the infrastructure level (a \texttt{GateStep} in a
  \texttt{PipelineSpec}, \S\ref{sec:pipelines}) is shown here to work at
  the application level too, and predates the infrastructure that now
  generalizes it.
\end{itemize}
\todo{Cite \texttt{agentic-ml-classification/docs/ml-classification.tex}
directly once it exists as a standalone artifact/preprint, rather than
inlining its content here.}

\subsection{Case Study II: Direct Agent-to-Agent Resource Scheduling}
\label{sec:casestudy2}
The resource-scheduler targets 6G-enabled factory-floor scheduling with six
agents: load monitor, task prioritization, resource allocation, failure
recovery, optimization, and human oversight. Its key architectural contrast
with Case Study I: it deliberately tests \emph{direct} agent-to-agent (A2A)
messaging for the negotiation-shaped parts of the workflow -- task
prioritization $\to$ resource allocation, and failure recovery $\to$
resource allocation, via a minimal in-process mailbox -- instead of routing
everything through a central orchestrator, while still keeping a
deterministic environment/state layer as the sole source of truth for
machine/slice state and constraint arithmetic. Every A2A send is still
logged as a traceable event; "direct" does not mean "unaudited": Resource
Allocation's constraint gate (no assignment to a machine under maintenance,
no slice over-capacity, checked cumulatively within a batch) applies
independently of what a sending agent claims, and Failure Recovery's
reroute proposals are re-validated against that exact same gate on receipt,
not accepted on the proposer's word.

\begin{itemize}
  \item \textbf{What this validates:} the propose/decide trust-boundary
  pattern is not tied to one specific orchestration topology -- it holds
  whether agents communicate through a central orchestrator (Case Study I)
  or negotiate peer-to-peer with a deterministic gate on receipt (Case
  Study II). This is directly relevant to the Sandbox's BYOA goal:
  partners' agents may not share one communication topology, and the
  Sandbox's trust model needs to accommodate both.
  \item \textbf{Status/scope caveat, verified 2026-09-01:} all six agents
  are individually implemented and unit-tested (92 deterministic tests,
  no live model call required, exercising only the environment/state
  layer) against synthetic and seeded-variance industrial scheduling data.
  A single-process orchestrator script now exists that chains all six
  agents through the shared A2A mailbox end-to-end, but as of this writing
  it has not yet been executed against a live model -- no run artifacts
  exist on disk -- so this case study is presented here as a preliminary,
  architectural data point validating the design pattern's applicability,
  not as full live-model evaluation evidence comparable to Case Study I's.
  \item \todo{A co-worker is currently porting resource-scheduler into
  \texttt{agent-sandbox} itself -- as a \texttt{gates/*.py} adapter module
  plus a \texttt{pipelines/*.yaml} spec, following the pattern documented
  in \texttt{agent-sandbox}'s README (already proven for
  \texttt{agentic\_ml}, see \texttt{docs/phase3-gates-handoff.md}) -- via
  the Sandbox UI rather than hand-written YAML. This work is local-only as
  of 2026-09-01, not yet pushed to the shared repository (per
  \texttt{docs/future-work-roadmap.md}'s Phase 4 rationale: \texttt{/agents}
  and \texttt{/pipelines} are global and shared, with no per-user isolation
  yet, so in-progress work stays local until it is ready to publish). Once
  that work lands, this subsection needs: (1) confirmation the ported
  pipeline runs end-to-end \emph{inside} the Sandbox itself, not just as
  standalone resource-scheduler scripts, and (2) whatever live-model
  evaluation results that testing produces, replacing this preliminary
  framing with real evaluation evidence.}
\end{itemize}

\subsection{What the Case Studies Show, Together}
Two independently-built systems, different domains, different orchestration
topologies, arrived at the same authority split. That is the argument for
why a shared Sandbox infrastructure -- rather than each partner re-deriving
this pattern independently -- has real value: the pattern is apparently not
domain-specific, which is exactly the premise Goal 2's holistic,
cross-workflow evaluation (\S\ref{sec:goals}) depends on.

% ============================================================================
\section{Intellectual Merit}
\label{sec:merit}
% ----------------------------------------------------------------------------

\subsection{Enhanced MCP Architecture}
MCP standardizes tool/data access but does not itself address
multi-institutional data governance. For Sandbox partners, participation is
contingent on preserving institutional data confidentiality, security, and
governance requirements that MCP alone does not enforce. We propose
extending the MCP architecture to support: (1) secure access to off-site and
distributed data stores; (2) institution-specific, policy-driven access
control, built on the same per-binding allowlist and credential-reference
mechanism already implemented (\S\ref{sec:infra}); and (3) mechanisms to
detect and prevent unintended data leakage across agents, tools, and
institutional boundaries.

\subsection{Multi-Agent Harness Engineering}
Harnesses mediate LLM$\leftrightarrow$tool/data/environment interaction to
improve reliability, but harness engineering today is largely manual and
requires deep domain
expertise~\cite{lin2026agenticharnessengineeringobservabilitydriven}. In MAS
specifically, an erroneous output from one agent can propagate to downstream
agents, causing cascading failures -- both case studies in
\S\ref{sec:casestudies} were hand-built to guard against exactly this, at
real engineering cost. Research direction: automated, domain-aware harness
engineering for MAS -- constraining agent actions, monitoring inter-agent
interactions, detecting erroneous behavior, and mitigating error propagation
across the workflow, ideally expressed generically at the
\texttt{PipelineSpec} level (\S\ref{sec:pipelines}) rather than re-derived
by hand for every new application.

\subsection{Emergent Cybersecurity Challenges in MAS}
MAS execute multi-step workflows, retain persistent state, and hold active
credentials -- expanding the attack surface beyond conventional software
vulnerabilities to include manipulation of reasoning processes, exploitation
of inter-agent communication protocols, and unauthorized or unintended tool
invocation. Existing cybersecurity frameworks, built for static,
human-directed software, do not adequately address this. \todo{Note
explicitly: this is proposed research, not yet evaluated in this paper --
the adversarial prompt-injection result in \S\ref{sec:casestudy1} is a
narrow, single-property test (state-bookkeeping integrity), not a general
MAS security evaluation. Don't conflate the two.}

\subsection{System-Level Innovation in MAS}
Even a modest set of agents admits an astronomically large space of possible
configurations and architectures -- one that grows far faster than the
capacity to evaluate it, and exhibits emergent behaviors (feedback loops,
bottlenecks, critical thresholds, phase transitions) that cannot be inferred
from individual agents' designs. This draws on systems-thinking literature
for rigorous system-level performance analysis, what-if evaluation, and
optimization -- Goal 2 (\S\ref{sec:goals}) made into a research problem, not
just a platform feature: the Sandbox's operator layer
(\S\ref{sec:design}) is the mechanism, but rigorous evaluation of the
resulting configuration space is itself open research.

\subsection{University-Industry Collaboration as an Innovation Lab}
The traditional tech-transfer/licensing model requires upfront agreement on
deliverables, IP sharing, and timelines before research happens -- often the
actual blocker to collaboration, not a shortage of good
ideas~\cite{hargadon2003breakthroughs}. The Sandbox's agentic-marketplace
model instead lets collaboration happen through experimentation itself:
partners showcase AI innovations as agents, connecting two-sided or
multi-sided based on complementary needs and competencies discovered
\emph{during} use, not negotiated in
advance~\cite{criado2018introduction,garciaarteagoitia2024experimentation}.

% ============================================================================
\section{Related Work}
\label{sec:related}
% ----------------------------------------------------------------------------
\todo{Needs an actual literature search -- the vision PDF's 58-entry
reference list is mostly surveys/position papers on agentic AI generally,
not competing \emph{systems}. The paragraphs below are a starting position,
not a completed survey.}

\textbf{Agent-environment/evaluation frameworks.} The Sandbox's
schema-first approach -- typed agent, run, and event specifications rather
than ad hoc harness code -- is the same design point independently
converged on by Meta's Agents Research Environments
(ARE)~\cite{froger2025arescalingagentenvironments}, which structures agent
evaluation around typed environments, apps, and events. The Sandbox targets
a different problem within that shared design pattern: multi-institutional,
multi-partner co-evaluation of competing agent implementations under shared
governance, rather than single-team agent benchmarking.

\textbf{Industry BYOA and low-code MAS platforms.} Two industry efforts
illustrate related but narrower design points. Langflow~\cite{langflow2026},
an open-source low-code visual builder for LLM applications, represents an
agent or RAG pipeline as a directed graph of components and exposes each
assembled flow as a callable API or MCP server; unlike the Sandbox, it
targets a single developer authoring one workflow rather than multiple
institutional partners co-evaluating competing agent implementations.
Plain's Bring-Your-Own-Agent (BYOA) feature~\cite{plain2026byoa} is
architecturally closer to the Sandbox's \texttt{SWAP\_AGENT} operator: it
decouples a customer-support platform's persistent infrastructure --
thread routing, escalation, SLAs -- from the AI agent handling
conversations, letting a team swap agents without rebuilding the
surrounding workflow. The Sandbox generalizes this swap principle beyond a
single vertical and single-tenant deployment to a multi-partner MAS setting
with typed event logging and institution-specific access control, a scope
neither Langflow nor Plain's BYOA was designed to address.

\textbf{MAS platforms and marketplaces.} Prior agent-marketplace and
agent-exchange proposals -- e.g.\ Chaffer et al.'s ``Hybrid Marketplace of
Ideas''~\cite{chaffer2025hybridmarketplaceideas} and Yang et al.'s ``Agent
Exchange''~\cite{yang2025agentexchangeshapingfuture} -- motivate the same
agentic-marketplace shift this paper argues for (\S\ref{sec:merit}).
\todo{Position the Sandbox's manufacturing/university-industry framing
against these more precisely once a fuller reading of both is done.}

\textbf{MAS in manufacturing.} Bio-inspired and self-organizing
manufacturing MAS predate LLM-driven
agents~\cite{Barbosa2010,Leitao2012,Bussmann2004}; recent LLM-agent work in
manufacturing~\cite{Farahani_2026} motivates the specific case study in
\S\ref{sec:casestudy1}. \todo{Reconcile this citation list with
\texttt{ml-classification.tex}'s own introduction, which cites the same
Farahani et al.\ work.}

\textbf{University-industry innovation lab models.} Non-AI-specific
experimentation-based collaboration
literature~\cite{criado2018introduction,garciaarteagoitia2024experimentation,hargadon2003breakthroughs}
motivates \S\ref{sec:merit}'s innovation-lab framing directly.

% ============================================================================
\section{Limitations and Future Work}
\label{sec:limitations}
% ----------------------------------------------------------------------------
\begin{itemize}
  \item \textbf{Infrastructure maturity.} As detailed in
  \S\ref{sec:status}, mid-run operation mutation, LLM-driven dynamic
  sub-agent spawning, sandboxed untrusted-gate execution, multi-tenant
  per-user isolation, and the cross-institution MCP governance gateway are
  not yet built. The paper's evidence for Sandbox \emph{value} currently
  rests substantially on the case studies (\S\ref{sec:casestudies})
  validating the underlying design pattern, not on the Sandbox
  infrastructure itself running multi-partner workloads yet -- be explicit
  about this gap rather than letting the two blur, even though the
  infrastructure that \emph{is} built (schemas, pipelines, operators,
  service/UI layer, 160 passing tests) is substantially more than a single
  headless runtime.
  \item \textbf{Resource-scheduler port and evaluation, in progress.} As
  detailed in \S\ref{sec:casestudy2}, a co-worker is actively porting
  resource-scheduler into \texttt{agent-sandbox} itself and testing it
  through the Sandbox UI; this work is local-only as of 2026-09-01. This is
  the single most important pending update to this draft before
  submission -- it will either strengthen Case Study II from an
  architectural data point into a real Sandbox-hosted evaluation, or
  surface gaps in the pipeline/gate mechanism (\S\ref{sec:pipelines}) that
  \S\ref{sec:infra} does not yet account for.
  \item \textbf{Case-study generalization.} Two case studies, one
  institution, one research team; no external partner has yet used the
  Sandbox pattern independently. Directly relevant since Goal 3/Goal 1
  (BYOA, marketplace) are unvalidated with a real second party.
  \item \textbf{Security.} \S\ref{sec:merit}'s cybersecurity agenda is
  proposed, not evaluated at the system level; the adversarial result in
  \S\ref{sec:casestudy1} tests one narrow property (state-bookkeeping
  integrity under prompt injection), not a general MAS threat model.
  \item \textbf{Remaining seed applications.} Predictive maintenance and
  planning/scheduling (\S\ref{sec:seedapps}) both now have case studies;
  supply chains and virtual organization does not, and is not yet built.
\end{itemize}

% ============================================================================
\section{Conclusion}
\label{sec:conclusion}
% ----------------------------------------------------------------------------
MAS adoption in manufacturing will depend on shared infrastructure for safe
experimentation as much as on any single application's performance. The
Sandbox proposes a specific, schema-first answer to what that
infrastructure looks like -- declarative agent and pipeline specs, a
harness-driven multi-agent runtime, a validated and audited pre-run operator
layer, and a structural isolation boundary -- verified today by 160 passing
automated tests, not merely designed on paper. Two independently-built
systems converging on its core trust-boundary pattern, without having
shared a codebase, is early but real evidence the idea holds past a single
implementation.

% ============================================================================
% \section*{Data Availability}
% \scaffold{Draft once decided: what's public vs. RIT-internal, dataset
% licensing (NGAFID/Kaggle), whether case-study code/run logs are
% shareable. Required-ish for SANER/FSE; good practice regardless.}
% ============================================================================

\bibliographystyle{IEEEtran}
\begin{thebibliography}{1}

\bibitem{visionpaper}
\todo{Cite the internal vision document (``AI Agentic Sandbox Project
Directions'') once it has a citable form, or fold its content in directly
without a self-citation if it stays internal.}

\bibitem{mlclassification}
\todo{Cite \texttt{agentic-ml-classification/docs/ml-classification.tex}
once that becomes a standalone preprint/paper.}

\bibitem{resourcescheduler}
\todo{Cite \texttt{resource-scheduler/README.md} / its own writeup once
Case Study II (\S\ref{sec:casestudy2}) is finalized with live-evaluation
results.}

\bibitem{Farahani_2026}
M.~A. Farahani, M.~I. Khan, T.~Wuest, Hybrid agentic AI and multi-agent
systems in smart manufacturing, Journal of Manufacturing Systems, vol.~86
(2026).

\bibitem{froger2025arescalingagentenvironments}
R.~Froger et al., ARE: Scaling Up Agent Environments and Evaluations,
arXiv:2509.17158 (2025).

\bibitem{nist_glossary_agent}
NIST, Glossary: Agent (AI). \url{https://csrc.nist.gov/glossary} (accessed
2026-09-01).

\bibitem{perilli2026agentic}
A.~Perilli et al., IDC FutureScape: Worldwide AI and Generative AI 2026
Predictions, IDC (2026).

\bibitem{langflow2026}
Langflow AI (DataStax), Langflow: An Open-Source Visual Framework for
Building Multi-agent and RAG Applications, GitHub repository (2026).

\bibitem{plain2026byoa}
S.~de Sousa, Introducing Bring Your Own Agent, Plain Blog, Feb.\ 2026.

\bibitem{chaffer2025hybridmarketplaceideas}
\todo{Full citation -- port from \texttt{docs/vision/references.bib}
(key: \texttt{chaffer2025hybridmarketplaceideas}).}

\bibitem{yang2025agentexchangeshapingfuture}
\todo{Full citation -- port from \texttt{docs/vision/references.bib}
(key: \texttt{yang2025agentexchangeshapingfuture}).}

\bibitem{Barbosa2010}
\todo{Full citation -- port from \texttt{docs/vision/references.bib}
(key: \texttt{Barbosa2010}).}

\bibitem{Leitao2012}
\todo{Full citation -- port from \texttt{docs/vision/references.bib}
(key: \texttt{Leitao2012}).}

\bibitem{Bussmann2004}
\todo{Full citation -- port from \texttt{docs/vision/references.bib}
(key: \texttt{Bussmann2004}).}

\bibitem{lin2026agenticharnessengineeringobservabilitydriven}
\todo{Full citation -- port from \texttt{docs/vision/references.bib}
(key: \texttt{lin2026agenticharnessengineeringobservabilitydriven}).}

\bibitem{hargadon2003breakthroughs}
\todo{Full citation -- port from \texttt{docs/vision/references.bib}
(key: \texttt{hargadon2003breakthroughs}).}

\bibitem{criado2018introduction}
\todo{Full citation -- port from \texttt{docs/vision/references.bib}
(key: \texttt{criado2018introduction}).}

\bibitem{garciaarteagoitia2024experimentation}
\todo{Full citation -- port from \texttt{docs/vision/references.bib}
(key: \texttt{garciaarteagoitia2024experimentation}).}

\todo{Port the remaining relevant entries from the vision PDF's 58-item
bibliography (systems-thinking literature, remaining MAS/manufacturing
citations) once Related Work (\S\ref{sec:related}) gets its real literature
search. Consider switching this document to \texttt{\textbackslash
bibliography\{../docs/vision/references\}} + BibTeX instead of hand-typed
\texttt{\textbackslash bibitem}s, now that most entries above are ported
straight from that shared \texttt{.bib} file anyway -- would remove the
duplication risk between the two.}

\end{thebibliography}

\end{document}
